\section{Results III: controlled data, failure modes and components}
\label{sec:controlled}

\paragraph{Where zero-shot forecasting wins and where it fails (RQ3).} Table~\ref{tab:tierc} gives, for each synthetic cell, the
geometric-mean MASE of each foundation model relative to Seasonal Naive, next to the best classical model of that cell. This is a
descriptive table: the classical model is chosen after seeing the results, which favours it. TimesFM~3 is better than the best of the
pre-specified classical baselines in 29 of the 38 cells (on the four intermittent-demand cells a zero forecast, which is not one of
them, does better still; Table~\ref{tab:zerod8}), often by a wide margin where classical methods lack information: a context of 256 hourly
points with weekly seasonality (0.371 against 0.914), a season longer than the context (0.519 against 0.956), two seasonalities
with one week of context (0.620 against 0.886), and a target driven by a known future covariate (0.336 against 0.752). On long
contexts and long horizons (up to 1{,}024 steps) the margin is a steady 6--8\%, and there TimesFM~3 and Chronos-2 are within 1\%
of each other.

The failures are specific; each comes from one synthetic cell and is descriptive, without a test. (i) \emph{Logistic growth} (D7). With 200 points of context the series has already
risen and starts to saturate; the true continuation is nearly flat, but TimesFM~3 forecasts a decline: its
median forecast over steps 25--48 lies below the last observation in 199 of the 200 replications, and its MASE over the 48 steps is 6.3 times that of
the naive forecast (TimesFM~2.5: 10.9; Chronos-2: 4.4; TiRex-2: 3.2; we did not inspect the direction of their errors). The simulation study we replicate reports the same spurious
decline \citep{ebrahim2026foundation}. With 48 points of context, which end before the steep part of the curve, TimesFM~3
under-forecasts the rise in almost every replication, and the best classical model (AutoETS) is about four times more accurate.
(ii) \emph{A regime change} near the origin, in one parametrisation (a level shift of five marginal standard deviations of the noise 24
steps before the origin, the same setting as the 5~sd, 24-step cell of the grid): TimesFM~3's relative MASE is 0.917 against 0.387 for both Chronos-2 and the best classical model, the
largest failure among the foundation models (TiRex-2 0.645, TimesFM~2.5 0.605; only Chronos-2 matches the best classical model),
although it still beats the seasonal naive forecast. With 24 points
after the shift the situation is ambiguous between a burst and a new level; a grid of shift sizes and timings of the same process in
Section~\ref{sec:robustness} shows that the weakness recurs across sizes and timings, not only in this cell. (iii) \emph{Level shifts} (D6, three of four cells) and the
\emph{random walk}: TimesFM~3 is 1--5\% worse than the best classical model. (iv) On \emph{white noise} all models match the historic mean, which is close to the optimal forecast (the true mean is known
only to the generator); no model does better. On intermittent demand (D8) the
verdict depends on the metric. Every foundation model has about a 20\% lower MASE than the naive forecast, the best classical
model on that metric. But demand is zero with probability 0.7, so the zero forecast is the true median, and it is at least as
good on MASE in every cell (relative to Seasonal Naive, 0.72--0.76 against 0.76--0.81 for the foundation models;
Table~\ref{tab:zerod8}, post hoc): the MASE advantage shows
that the models forecast a median near zero, not that they forecast demand well. On the root mean squared scaled error, which
rewards the conditional mean rather than the median, the historic mean and AutoETS are better than TimesFM~3 (0.709 against 0.811
with 200 points and $H = 48$). On the quantile loss the foundation models are better than the probabilistic classical baselines
we ran (0.567 against 0.674), none of which is designed for intermittent demand. Their quantiles therefore look promising against
these baselines, while the returned median is not an appropriate forecast for decisions governed by squared error (analysis
added after the freeze, not pre-specified).

\begin{table}[!htbp]
\centering
\caption{Tier C (descriptive): geometric-mean MASE relative to Seasonal Naive over 200 replications per cell, and the best of the
pre-specified classical baselines in each cell, chosen after the fact (on D8 the zero forecast is better; Table~\ref{tab:zerod8}). Last column: TimesFM~3 over the best classical model (bold: TimesFM~3
worse). TiRex-2 is not run beyond 320 steps.}
\label{tab:tierc}
\scriptsize
\setlength{\tabcolsep}{3pt}
\begin{tabular}{lrrrrrlr}
\toprule
Cell & $H$ & TimesFM-3 & Chronos-2 & TiRex-2 & TimesFM-2.5 & Best classical & TFM-3 / best \\
\midrule
\multicolumn{8}{l}{\textit{C-horizon}} \\
uni long\_horizon 1024 & 1024 & 0.240 & 0.241 & -- & 0.252 & 0.257 (MSTL-ARIMA) & 0.937 \\
uni long\_horizon 256 & 256 & 0.205 & 0.205 & 0.218 & 0.208 & 0.222 (MSTL-ARIMA) & 0.926 \\
uni long\_horizon 512 & 512 & 0.239 & 0.239 & -- & 0.243 & 0.257 (MSTL-ARIMA) & 0.928 \\
\midrule
\multicolumn{8}{l}{\textit{C-ladder}} \\
uni context\_ladder 1024 & 168 & 0.328 & 0.329 & 0.353 & 0.334 & 0.356 (MSTL-ARIMA) & 0.919 \\
uni context\_ladder 15360 & 168 & 0.252 & 0.255 & 0.266 & 0.255 & 0.275 (MSTL-ARIMA) & 0.919 \\
uni context\_ladder 256 & 168 & 0.371 & 0.350 & 0.744 & 0.357 & 0.914 (MSTL-ETS) & 0.405 \\
uni context\_ladder 4096 & 168 & 0.232 & 0.232 & 0.243 & 0.235 & 0.251 (MSTL-ARIMA) & 0.924 \\
uni context\_ladder 64 & 168 & 0.883 & 0.994 & 0.914 & 0.801 & 0.960 (MSTL-ETS) & 0.919 \\
\midrule
\multicolumn{8}{l}{\textit{C-mv}} \\
mv independent 8 & 24 & 0.783 & 0.785 & 0.783 & 0.786 & 0.787 (AutoARIMA) & 0.995 \\
mv known\_cov snr0.5 & 24 & 0.684 & 0.705 & 0.686 & 0.766 & 0.764 (AutoARIMA) & 0.896 \\
mv known\_cov snr2 & 24 & 0.336 & 0.338 & 0.343 & 0.754 & 0.752 (Mean) & 0.447 \\
mv lead\_lag 1 & 24 & 0.843 & 0.850 & 0.846 & 0.853 & 0.872 (AutoARIMA) & 0.967 \\
mv lead\_lag 24 & 24 & 0.677 & 0.572 & 0.848 & 0.861 & 0.885 (AutoARIMA) & 0.764 \\
mv lead\_lag 5 & 24 & 0.807 & 0.767 & 0.828 & 0.841 & 0.889 (AutoARIMA) & 0.908 \\
mv spurious\_past\_cov & 24 & 0.747 & 0.750 & 0.747 & 0.751 & 0.751 (AutoARIMA) & 0.995 \\
mv var1 2 & 24 & 0.739 & 0.739 & 0.744 & 0.741 & 0.741 (AutoARIMA) & 0.997 \\
mv var1 32 & 24 & 0.740 & 0.738 & 0.742 & 0.743 & 0.743 (Mean) & 0.997 \\
mv var1 64 & 24 & 0.744 & 0.744 & 0.746 & 0.746 & 0.745 (Mean) & 0.999 \\
mv var1 8 & 24 & 0.755 & 0.751 & 0.758 & 0.759 & 0.758 (AutoARIMA) & 0.995 \\
\midrule
\multicolumn{8}{l}{\textit{C-replication}} \\
rep D6 n200 & 48 & 0.902 & 0.927 & 1.178 & 0.871 & 0.895 (AutoETS) & \textbf{1.008} \\
rep D6 n200 h12 & 12 & 0.912 & 0.907 & 0.906 & 0.898 & 0.894 (AutoETS) & \textbf{1.020} \\
rep D6 n48 & 48 & 1.000 & 1.186 & 1.148 & 0.977 & 0.958 (AutoETS) & \textbf{1.044} \\
rep D6 n48 h12 & 12 & 0.886 & 0.891 & 0.853 & 0.860 & 0.931 (AutoETS) & 0.952 \\
rep D7 n200 & 48 & 6.322 & 4.382 & 3.245 & 10.924 & 1.000 (Naive) & \textbf{6.322} \\
rep D7 n200 h12 & 12 & 1.024 & 1.136 & 0.993 & 1.257 & 1.000 (Naive) & \textbf{1.024} \\
rep D7 n48 & 48 & 0.714 & 0.458 & 0.712 & 0.876 & 0.183 (AutoETS) & \textbf{3.901} \\
rep D7 n48 h12 & 12 & 0.170 & 0.393 & 0.246 & 0.307 & 0.178 (AutoTheta) & 0.950 \\
rep D8 n200 & 48 & 0.767 & 0.768 & 0.775 & 0.767 & 1.000 (Naive) & 0.767 \\
rep D8 n200 h12 & 12 & 0.773 & 0.774 & 0.787 & 0.769 & 1.000 (Naive) & 0.773 \\
rep D8 n48 & 48 & 0.807 & 0.783 & 0.796 & 0.795 & 1.000 (Naive) & 0.807 \\
rep D8 n48 h12 & 12 & 0.768 & 0.759 & 0.780 & 0.762 & 1.000 (Naive) & 0.768 \\
\midrule
\multicolumn{8}{l}{\textit{C-uni}} \\
uni long\_period ctx200 & 60 & 0.519 & 0.509 & 0.974 & 0.731 & 0.956 (AutoTheta) & 0.543 \\
uni outliers & 48 & 0.278 & 0.282 & 0.305 & 0.284 & 0.326 (MSTL-ARIMA) & 0.853 \\
uni random\_walk & 64 & 1.051 & 1.021 & 1.042 & 1.018 & 1.000 (Naive) & \textbf{1.051} \\
uni regime\_change & 48 & 0.917 & 0.387 & 0.645 & 0.605 & 0.387 (MSTL-Theta) & \textbf{2.368} \\
uni two\_seasons ctx168 & 168 & 0.620 & 0.597 & 0.847 & 0.530 & 0.886 (MSTL-ARIMA) & 0.700 \\
uni two\_seasons ctx672 & 168 & 0.315 & 0.313 & 0.338 & 0.322 & 0.346 (MSTL-ARIMA) & 0.911 \\
uni white\_noise & 64 & 0.714 & 0.714 & 0.714 & 0.714 & 0.713 (Mean) & \textbf{1.001} \\
\bottomrule
\end{tabular}

\end{table}

\begin{table}[!htbp]
\centering
\caption{Intermittent demand (D8, descriptive, not pre-specified): geometric-mean MASE relative to Seasonal Naive (the naive forecast
here) of the zero forecast, which is the true median (demand is zero with probability 0.7), and of the foundation models and the
historic mean, on the replications with a positive error for every model.}
\label{tab:zerod8}
\scriptsize
\begin{tabular}{lrrrrrrr}
\toprule
Cell & Reps & Zero & TimesFM-3 & Chronos-2 & TiRex-2 & TimesFM-2.5 & Mean \\
\midrule
rep D8 n200 & 200 & 0.757 & 0.767 & 0.768 & 0.775 & 0.767 & 1.079 \\
rep D8 n200 h12 & 197 & 0.759 & 0.773 & 0.774 & 0.787 & 0.769 & 1.138 \\
rep D8 n48 & 200 & 0.762 & 0.807 & 0.783 & 0.796 & 0.795 & 1.099 \\
rep D8 n48 h12 & 197 & 0.722 & 0.768 & 0.759 & 0.780 & 0.762 & 1.159 \\
\bottomrule
\end{tabular}

\end{table}

\paragraph{H3: the detrending stage.} The pre-specified test (\texttt{tfm3-nodetrend} against \texttt{tfm3}, D7, 200 points,
$H = 48$, steps 25--48) gives a ratio of exactly \Hthree{}: the two configurations produce identical forecasts on all 200
replications, so H3 is uninformative (formally not supported). The reason is a design error of the protocol. The detrending gate opens only when a line
explains more than 75\% of the context's variance, and in this cell the 200-point context is a long flat start followed by the
rise, which a line fits with $R^2 < 0.75$ in all 200 contexts. The gate opens in every context of the other three D7 cells. There,
in an exploratory analysis without multiplicity adjustment, switching the detrending off lowers MASE over steps 25--48 by 13\%
(0.873, 0.810--0.941) and over the whole 48-step horizon by 17\% (0.830) with 48 points of context, and by 7\% at $H = 12$ with 200
points (0.930), but almost doubles it at $H = 12$ with 48 points (1.888). Across these configurations, disabling the detrending helps in some settings and hurts in others; the design does not isolate
horizon length from the position of the logistic transition, whose inflection point is at $0.6(n + H)$, so a longer horizon also
changes the observed part of the curve. We did not analyse the forecasts further. It is not the cause of the largest D7 failure: in the 200-point cell, where TimesFM~3 forecasts
the spurious decline, the gate is closed and the forecast comes from the network alone.

\paragraph{Multivariate mode (H4, RQ5).} Joint forecasting helps when there is something to learn across variates: for the
follower of a lead-lag pair with lag 24, MASE falls to \Hfoura{} (\HfouraLo--\HfouraHi) of the univariate forecast, and with a known
future covariate (signal-to-noise 2) to \Hfourb{} (\HfourbLo--\HfourbHi). In two controls that add no predictive information it costs no accuracy
beyond the $\pm$3\% margin: with eight independent series stacked as variates the ratio is \Hfourc{}, and with a redundant noisy
past covariate (a noisy copy of the target over the context, called spurious in the protocol) \Hfourd{}, both equivalent within $\pm$3\% (Table~\ref{tab:familyh}); the result is limited to these two kinds of uninformative input; other processes without a signal were not tested. On a lead-lag pair with lag 24, Chronos-2 reaches a lower MASE than
TimesFM~3 (Table~\ref{tab:tierc}); we did not measure its multivariate gain separately. On real data the gain is about 3\% lower SQL on the group tasks (P4a), a decision that does not survive the small-cluster and
non-overlapping-window checks, and no gain from oracle covariates is resolved (S3; the interval includes both gains and losses);
their relation to the target is weaker and less direct than in the synthetic cells.

\begin{table}[!htbp]
\centering
\caption{Family H: primary tests of the pre-declared hypotheses (Holm over seven tests). H4c and H4d are equivalence tests with a
$\pm$3\% margin (two one-sided tests): for them the table gives the 90\% interval (marked $^{90}$), and $p_{\mathrm{Holm}}$ adjusts the TOST $p$; H5 is a
non-superiority test with floor 0.90 and a one-sided $p$. H3 is uninformative: its cell never activates the mechanism it tests
(ratio exactly 1). Ratios to four decimals.}
\label{tab:familyh}
\scriptsize
\setlength{\tabcolsep}{3pt}
\begin{tabular}{lllrlll}
\toprule
 & Contrast & Prediction & Ratio & 95\% CI & $p_{\mathrm{Holm}}$ & Decision \\
\midrule
H2 & W-seen vs W-post DiD, TimesFM-3 vs Chronos-2 & ratio $<$ 1 & 0.9799 & [0.9362, 1.0207] & 0.714 & not supported \\
H3 & no detrending vs default, D7, steps 25--48 & ratio $<$ 1 & 1.0000 & [1.0000, 1.0000] & 1.000 & uninformative \\
H4a & multivariate vs univariate, lead-lag 24 & ratio $<$ 1 & 0.6159 & [0.5913, 0.6425] & $<$0.001 & supported \\
H4b & multivariate vs univariate, known covariate & ratio $<$ 1 & 0.4458 & [0.4315, 0.4604] & $<$0.001 & supported \\
H4c & multivariate vs univariate, independent series & within $\pm$3\% & 0.9979 & [0.9968, 0.9990]$^{90}$ & $<$0.001 & supported \\
H4d & multivariate vs univariate, spurious covariate & within $\pm$3\% & 1.0004 & [0.9999, 1.0009]$^{90}$ & $<$0.001 & supported \\
H5 & TimesFM-3 uni vs Naive, FX and FRED & not $<$ 0.90 & 1.1096 & [1.0489, 1.1784] & $<$0.001 & supported \\
\bottomrule
\end{tabular}

\end{table}

\subsection{Ablations (RQ6)}
\label{sec:ablations}

Each ablation changes one fixed step or option of \texttt{tfm3} and is compared with it on the tasks where the change can act
(Table~\ref{tab:ablations}). Apart from the variate attention, discussed at the end, the effects are small. On the synthetic data, switching off the detrending raises MASE by 0.7\% and
padding future covariates by repetition (\texttt{tfm3-edge}) by 0.7\%; clipping at zero lowers it by 0.05\%; all three differences
are resolved but too small to matter in practice. Switching off the CPM refinement or the stitching, or adding symmetric averaging or
z-normalisation, changes MASE by less than 0.3\% and none of these differences is resolved. On the 20\% sample of Tier B series no
ablation has a resolved effect. Averages hide where a step acts (descriptive, not pre-specified). We measure a change as the largest absolute difference between
the median forecasts, divided by the MASE scale of the context; numerical noise, measured the same way between two runs of the
identical computation (the invalid ablation of deviation 2), is zero on most windows but reaches 1.9\% of the scale on Tier B
($9 \cdot 10^{-6}$ on Tier C), and a window counts as changed only above that maximum. The detrending gate, computed from each context with
the model's own rule, opens on 2.8\% of the Tier B windows and 8.9\% of the synthetic ones, and the forecasts change on the same
windows (plus four on Tier B). Where it opens, switching it off lowers MASE by 12.8\% on Tier B, raising it in only 21\% of those
windows, and raises MASE by 8.1\% on Tier C, in 47\% of them: on the real windows where it acts the fixed linear extrapolation
mostly hurts, on the synthetic ones it helps on average but in barely half of the windows. Clipping at zero changes 0.6\% of the
Tier B and 4.2\% of the synthetic forecasts (on the latter it lowers MASE by 1.1\%). Symmetric averaging, stitching and covariate
edge handling change (nearly) every forecast they can act on, CPM and iterative RevIN 14\% of the Tier B ones beyond the noise,
and their effects stay below 0.7\% in absolute value as the geometric mean, with equal weight per window, of the ratio of MASE to
that of \texttt{tfm3} over the windows they change; z-normalisation changes 0.9\% of the Tier B forecasts and 12\% of the synthetic ones,
without effect on the latter; under the pre-specified estimator of Table~\ref{tab:ablations} the estimates reach
$-3.4\%$ for covariate edge handling and $-1.2\%$ for the detrending on Tier B, neither resolved. The fixed steps of the inference path therefore matter in specific situations (the detrending gate,
H3) but, among the tasks and losses analysed, change average accuracy little, partly because they act on few windows and their
effects there differ by setting; an unresolved difference is not shown to be
small, and equivalence was not tested. A reading, not a result: the network does most of the work. Switching off the
variate attention is different in kind (re-run after a software error, Appendix~\ref{app:deviations}): it removes a trained sublayer
from all 20 layers, and the forecasts become unusable, with MASE 795 times that of \texttt{tfm3} on Tier C and 10{,}881
times on Tier B. Even for a single series the sublayer is part of the network (attention over one variate returns that
variate's value projection, which is added to the residual stream), so this ablation measures the network's dependence on the
sublayer, not the value of information from other series; that value is measured by the multivariate against univariate contrasts
(H4, P4a).

\begin{table}[!htbp]
\centering
\caption{Ablations (secondary family S5, Holm over the 16 tests of both tiers; ratios to four decimals; with Holm within each tier,
\texttt{tfm3-sym} on Tier C, unadjusted $p = 0.006$, Holm-adjusted 0.023, would also be resolved): ratio of MASE, ablated configuration / \texttt{tfm3}, on the
Tier B sample and on Tier C, on the tasks where the change can act.}
\label{tab:ablations}
\scriptsize
\setlength{\tabcolsep}{4pt}
\begin{tabular}{llrllrll}
\toprule
Configuration & Change & Ratio B & 95\% CI & $p_{\mathrm{Holm}}$ & Ratio C & 95\% CI & $p_{\mathrm{Holm}}$ \\
\midrule
\texttt{tfm3-nodetrend} & linear detrending off & 0.9885 & [0.9687, 1.0001] & 0.895 & 1.0069 & [1.0036, 1.0102] & 0.002 \\
\texttt{tfm3-nocpmrevin} & CPM and iterative RevIN off & 0.9987 & [0.9946, 1.0020] & 1.000 & 0.9999 & [0.9984, 1.0013] & 1.000 \\
\texttt{tfm3-novarattn} & variate attention off & 10{,}881 & [8{,}595, 13{,}682] & 0.002 & 795 & [784, 806] & 0.002 \\
\texttt{tfm3-nostitch} & stitching off & 1.0063 & [0.9962, 1.0160] & 1.000 & 1.0004 & [1.0001, 1.0007] & 0.202 \\
\texttt{tfm3-sym} & symmetric averaging on & 0.9995 & [0.9708, 1.0300] & 1.000 & 1.0023 & [1.0006, 1.0040] & 0.064 \\
\texttt{tfm3-pos} & make\_positive on & 1.0000 & [0.9999, 1.0002] & 1.000 & 0.9995 & [0.9989, 0.9999] & 0.002 \\
\texttt{tfm3-znorm} & z-normalised inputs & 1.0006 & [0.9995, 1.0027] & 1.000 & 1.0000 & [1.0000, 1.0000] & 1.000 \\
\texttt{tfm3-edge} & covariate edge handling & 0.9656 & [0.8951, 1.0380] & 1.000 & 1.0069 & [1.0034, 1.0106] & 0.002 \\
\bottomrule
\end{tabular}

\end{table}
